\chapter{Conclusions}

This thesis has investigated the origin of the persistent decline of the inner density observed in collisionless dark matter halos evolved in isolation for 100 Gyr using the SWIFT code. The central conclusion of this work is that the continuous reduction of the inner density in these simulated halos is not a physical phenomenon, nor is it the result of post-processing errors or algorithmic biases within the gravity solver. Across the full set of tests, the core continues to evolve after its initial formation rather than settling into a stable inner configuration. This behaviour remains visible when the box size is changed, when the initial halo structure is modified, when the gravitational softening length is varied, and when the self-gravity solver parameters are tightened or relaxed. In addition, the rise in central velocity dispersion is consistent with numericalheating rather than adiabatic expansion.

In summary, the persistent inner decline in the halo is most naturally explained by residual numerical heating, specifically finite-$N$ two-body relaxation and discreteness noise within the innermost regions of the halo. Both theory and simulations indicate that finite-$N$ two-body relaxation should secularly flatten cuspy profiles unless $N$ is astronomically large or $\varepsilon$ is very coarse. Notably, removing long range forces had only a small impact until scales $\sim \varepsilon$ were reached.

\section{Future work}
Building upon these findings, several avenues for future research can be pursued. Since the observed numerical heating is fundamentally a finite-$N$ effect, future studies should execute very high resolution simulations with massively increased particle counts. This approach is necessary to empirically determine the threshold at which two-body relaxation becomes negligible over both cosmological and extended timescales.

Additionally, investigating the implementation of adaptive gravitational softening or alternative softening kernel shapes within SWIFT could provide insights into better suppressing short range noise without compromising the accuracy of the global potential.

A more immediate and practical improvement is related to the hybrid model simulations used in this study. Future work should focus on finding a more systematic method to obtain a smoother potential distribution when coupling the analytical external potential with the inner $N$-body system. Currently, achieving a smooth transition relies heavily on trial-and-error tinkering with the parameters of the external potential module.

Finally, future work should introduce gas dynamics, star formation, and stellar feedback into the simulations. Baryonic processes can lead to physical core formation (e.g., via rapid supernova feedback) or central contraction. Understanding exactly how baseline numerical heating interacts with these physical baryonic effects is crucial for the reliability of modern cosmological hydrodynamical simulations.